\section{Implementación}
\label{sec:implementacion}

El sistema está escrito en Java~17 sin dependencias externas. El código
fuente está en \texttt{src/triage/}, las pruebas en \texttt{test/triage/} y
el benchmark en \texttt{bench/triage/}; el \texttt{README.md} del
repositorio contiene los comandos exactos de compilación y ejecución. La
Tabla~\ref{tab:estado} resume el estado de cada funcionalidad.

\begin{table}[t]
  \caption{Estado de las funcionalidades en la segunda entrega}
  \label{tab:estado}
  \centering
  \footnotesize
  \begin{tabular}{p{0.52\columnwidth}p{0.18\columnwidth}p{0.16\columnwidth}}
    \toprule
    Funcionalidad & Entrega 1 & Entrega 2 \\
    \midrule
    Registrar solicitud (RF1) & Sí & Sí \\
    Atender por prioridad (RF2) & Sí & Sí \\
    Buscar por identificador (RF3) & Sí & Sí \\
    Cancelar solicitud (RF4) & Sí & Sí \\
    Deshacer última atención (RF5) & Pendiente & Sí \\
    Consultar historial (RF6) & Sí & Sí \\
    Listar pendientes en orden (RF7) & --- & Sí \\
    Índice por identificador en $O(1)$ & Pendiente & Sí \\
    Validación de entradas & Parcial & Sí \\
    Pruebas automatizadas & Parcial & Sí \\
    Análisis experimental & Pendiente & Preliminar \\
    \bottomrule
  \end{tabular}
\end{table}

\textbf{Implementado.} Los siete requerimientos funcionales están
disponibles en \texttt{TriageSystem}: \texttt{registerRequest},
\texttt{serveNext}, \texttt{peekNext}, \texttt{findById},
\texttt{cancelRequest}, \texttt{undoLastService}, \texttt{getHistory} y
\texttt{getPendingInOrder}. El menú de \texttt{Main} expone todas estas
operaciones. Las estructuras de soporte (\texttt{LinkedQueue},
\texttt{LevelPriorityQueue}, \texttt{HashTable} y \texttt{LinkedStack}) son
de implementación propia; de la biblioteca estándar solo se usa
\texttt{ArrayList} para devolver listados. El Listado~\ref{lst:undo}
muestra la operación de deshacer.

\begin{lstlisting}[caption={Deshacer en \texttt{TriageSystem} (mensaje abreviado)},label={lst:undo},float=t]
public Ticket undoLastService() {
  Ticket last = history.lastServed();
  if (last == null) {
    return null;
  }
  String id = last.getId();
  if (index.containsKey(id)) {
    throw new IllegalStateException(
        "No se puede deshacer: " + id);
  }
  history.removeLast();
  index.put(id, queue.insertFront(last));
  return last;
}
\end{lstlisting}

\textbf{Validación de entradas.} Se rechazan, con una excepción
\texttt{IllegalArgument\-Exception}, los identificadores y nombres nulos,
vacíos o en blanco, las urgencias fuera de 1--5 y los identificadores
repetidos entre pendientes; el menú además captura urgencias no numéricas.
Los identificadores se normalizan quitando espacios iniciales y finales,
para que ``\texttt{C1}'' y ``\texttt{\ C1\ }'' no se registren como distintos.

\textbf{Decisiones técnicas relevantes.}
\begin{itemize}
  \item \emph{El índice vive en la fachada}, no dentro de la cola de
    prioridad: la cola solo sabe ordenar y la fachada es la única que
    coordina cola, índice e historial, lo que deja el invariante en un solo
    lugar. \texttt{TriageSystem} incluye un método de verificación del
    invariante que usan las pruebas.
  \item \emph{Localizadores opacos y verificados.} El nodo expuesto por
    \texttt{LinkedQueue} solo permite leer su valor; retirar un nodo que no
    pertenece a la cola (o que ya fue retirado) lanza una excepción.
  \item \emph{Conflicto al deshacer.} Se eligió rechazar la operación en
    lugar de sobrescribir la solicitud pendiente con el mismo
    identificador, porque sobrescribir haría perder una solicitud de forma
    silenciosa.
  \item \emph{Búsqueda lineal conservada como línea base.} La búsqueda
    de la Entrega 1 (\texttt{linearFindById} en
    \texttt{LevelPriorityQueue}) ya no la usa el sistema, pero se conserva para el benchmark y para
    contrastar el índice en las pruebas.
  \item \emph{Pruebas sin JUnit.} Se escribió un mini marco de aserciones
    (\texttt{Check.java}) para que el proyecto compile y se pruebe solo con
    el JDK, sin descargar dependencias.
\end{itemize}

\textbf{Parcial o pendiente.} No queda ningún requerimiento funcional sin
implementar. El análisis experimental es preliminar: se ejecutó en un solo
equipo y con datos de urgencia uniforme (Sección~\ref{sec:resultados}).
Las funcionalidades previstas para la entrega final se describen en la
Sección~\ref{sec:estado}.
